\chapter{A Concrete Naming Certificate for $\CompactIntervalModulusUp$}
\label{ch:concrete-instances-compact-interval-modulus-namecert}
\origin{human}

Following Bishop and Bridges, the $\CompactIntervalModulusUp$ packet records the finite modulus route for compact intervals as displayed NameCert data. It refines the compact interval source of \autoref{ch:concrete-instances-bishop-compact-interval-namecert}, consumes the finite compactness surface of \autoref{ch:concrete-instances-closed-bounded-interval-compact-namecert}, and terminates through the real-name seal of \autoref{ch:concrete-instances-real-namecert}. The packet supplies a choice-free L10 route from interval compactness to modulus readback for downstream Completion, LocatedMetric, and separated-reflection consumers without asserting open-cover compactness, selected subcovers, quotient real equality, countable choice, or ambient $\RealUp$ completeness.

\begin{definition}[Compact interval modulus carrier]
\label{def:compact-interval-modulus-carrier}
A $\CompactIntervalModulusUp$ carrier is a finite $\BHist$ packet
$$
M=(I,K,F,T,W,R,E,H,C,P,N)
$$
with the following displayed rows:
\begin{enumerate}
\item $I$ is the $\BishopCompactIntervalUp$ source row, exposing the constructive interval endpoint, locatedness, finite-net, dyadic, stream, regular, and Real-facing data of the compact interval route.
\item $K$ is the $\ClosedBoundedIntervalCompactUp$ compactness row over the same interval source.
\item $F$ is the finite-net handoff row, recording the inspected mesh, radius ledger, and coverage comparison used by modulus consumers.
\item $T$ is the dyadic tolerance row for requested output precision.
\item $W$ is the $\StreamNameUp$ window row opened by the interval and finite-net data.
\item $R$ is the $\RegSeqRatUp$ readback row attached to the finite windows and tolerance ledger.
\item $E$ is the terminal $\RealUp$ seal for endpoint and modulus readback.
\item $H$ records componentwise $\hsame$ transport for source interval, compactness, finite-net, tolerance, stream, regular, Real, replay, provenance, and naming rows.
\item $C$ records displayed $\Cont$ replay from interval inspection through compact finite-net handoff, tolerance request, finite windows, regular readback, and Real sealing.
\item $P$ is the $\Pkg$ provenance over $\CompactIntervalModulusUp$, $\BishopCompactIntervalUp$, $\ClosedBoundedIntervalCompactUp$, $\DyadicRatCoreUp$, $\StreamNameUp$, $\RegSeqRatUp$, $\RealUp$, $\BHist$, $\hsame$, $\Cont$, and $\NameCert$ rows.
\item $N$ is the local $\NameCert_{\CompactIntervalModulusUp}$ row.
\end{enumerate}
The classifier accepts only packets with these displayed rows. It has no coordinate for arbitrary open covers, selected finite subcovers, quotient real equality, countable-choice schedules, selected representatives, or ambient $\RealUp$ completeness.
\end{definition}

\begin{theorem}[Compact interval modulus NameCert obligations]
\label{thm:compact-interval-modulus-namecert-obligations}
For every accepted $\CompactIntervalModulusUp$ carrier $M=(I,K,F,T,W,R,E,H,C,P,N)$, the local $\NameCert_{\CompactIntervalModulusUp}$ surface consists exactly of carrier admission, Bishop compact interval source exposure through $I$, closed-bounded compactness exposure through $K$, finite-net handoff through $F$, dyadic tolerance control through $T$, finite window exposure through $W$, regular rational readback through $R$, terminal Real sealing through $E$, componentwise transport through $H$, displayed replay through $C$, provenance through $P$, and local naming through $N$.
\end{theorem}
\begin{proof}
Carrier admission is projection from \autoref{def:compact-interval-modulus-carrier}. The tuple lists the interval source, compactness row, finite-net handoff, tolerance row, stream windows, regular readback, Real seal, and structural rows.

The interval source is read first through $I$ and $K$. These rows bind the Bishop compact interval route to the closed-bounded compactness surface before a modulus consumer can inspect any mesh or tolerance data.

The modulus handoff is finite. The row $F$ records the mesh and coverage comparison, while $T,W,R$ route each precision request through dyadic tolerance, finite stream windows, and regular rational readback.

Finally $E,H,C,P,N$ seal, transport, replay, source, and name the same route. Since the carrier has no open-cover, selected-subcover, quotient, choice, representative, or ambient-completeness coordinate, the listed rows exhaust the local certificate surface.
\end{proof}

\begin{theorem}[Compact interval modulus finite-net handoff]
\label{thm:compact-interval-modulus-finite-net-handoff}
Every modulus-facing consumer of an accepted $\CompactIntervalModulusUp$ packet factors through the finite route
$$
\BishopCompactIntervalUp
\longmapsto
\ClosedBoundedIntervalCompactUp
\longmapsto
\DyadicRatCoreUp
\longmapsto
\StreamNameUp
\longmapsto
\RegSeqRatUp
\longmapsto
\RealUp .
$$
The handoff exposes only the compact interval source, closed-bounded compactness row, finite-net mesh, dyadic tolerance, finite windows, regular readback, Real sealing, transport, replay, provenance, and local naming; it exports no open-cover compactness theorem, selected subcover, quotient real equality, countable-choice schedule, selected representative, or ambient Real completeness principle.
\end{theorem}
\begin{proof}
Project the accepted carrier $M=(I,K,F,T,W,R,E,H,C,P,N)$. The compact interval source is $I$, the compactness row is $K$, and the finite-net handoff is $F$.

Read a modulus request through $T,W,R$. The dyadic tolerance row fixes the requested precision, the stream row exposes the finite windows permitted by the interval and net data, and the regular row gives the rational readback.

The terminal route reaches $E$ only after those displayed rows have been consumed. The structural rows $H,C,P,N$ preserve the same order for downstream consumers.

Because the carrier lacks open-cover, selected-subcover, quotient, choice, representative, and ambient-completeness coordinates, no stronger compactness or Real principle can be read from the handoff.
\end{proof}

\begin{theorem}[Compact interval modulus completion readback]
\label{thm:compact-interval-modulus-completion-readback}
Any $\CompletionUp$ or $\LocatedMetricUp$ consumer reading a modulus from an accepted $\CompactIntervalModulusUp$ packet receives the displayed rows $I,K,F,T,W,R,E,H,C,P,N$ and no additional equality, selection, or completeness data. A $\SeparatedMetricUp$ consumer may later transport the distance or readback rows through its own reflection surface, but this packet itself does not turn compact interval modulus evidence into host equality, quotient equality, selected representatives, selected subcovers, countable choice, or ambient $\RealUp$ completeness.
\end{theorem}
\begin{proof}
Start with the carrier of \autoref{def:compact-interval-modulus-carrier}. Its completion-facing data are exactly the compact interval source, compactness row, finite-net handoff, tolerance row, stream row, regular readback, Real seal, and structural rows.

Apply \autoref{thm:compact-interval-modulus-namecert-obligations}. The obligation surface forces source exposure through $I,K$, handoff through $F$, tolerance and readback through $T,W,R$, and terminal sealing through $E$.

Apply \autoref{thm:compact-interval-modulus-finite-net-handoff}. The theorem orders the same rows as the finite route from compact interval data to modulus readback.

Thus Completion and LocatedMetric consumers may consume the displayed distance and readback data only through this route. A later SeparatedMetric consumer must add its own reflection rows before any identity interpretation is available, and all refused open-cover, selection, quotient, choice, representative, and ambient-completeness objects remain absent.
\end{proof}

\falsifiablePrediction{Every accepted $\CompactIntervalModulusUp$ read passes through BishopCompactIntervalUp, ClosedBoundedIntervalCompactUp, DyadicRatCoreUp, StreamNameUp, RegSeqRatUp, RealUp, transport, replay, provenance, and local naming rows before any modulus output is visible.}

\independenceWitness{bishop-compact-interval, closed-bounded-interval-compact, compact-interval-modulus-bridge, compact-interval-choice-free, uniformmodulus: $\CompactIntervalModulusUp$ is not bijective to any listed sibling because it jointly binds compact interval source rows, closed-bounded compactness rows, finite-net handoff, dyadic tolerance, stream windows, regular readback, terminal Real sealing, transport, replay, provenance, and local naming.}

\closureat{CompactIntervalModulusUp}{seedStr}

\begin{closurestatus}{\CompactIntervalModulusUp}
  \constructivestory{A $\CompactIntervalModulusUp$ packet is generated from BishopCompactIntervalUp, ClosedBoundedIntervalCompactUp, finite-net handoff, DyadicRatCoreUp tolerance, StreamName windows, RegSeqRat readback, a RealUp seal, componentwise hsame transport, displayed Cont replay, Pkg provenance, and a local NameCert row.}
  \theoryclosure{\seedClosure}
  \scopeclosed{\autoref{def:compact-interval-modulus-carrier}, \autoref{thm:compact-interval-modulus-namecert-obligations}, \autoref{thm:compact-interval-modulus-finite-net-handoff}, and \autoref{thm:compact-interval-modulus-completion-readback} seed the finite compact-interval modulus route over compact interval, closed-bounded compactness, dyadic, stream, regular, and Real rows.}
  \formalstatus{\unformalizedV}
  \bridgestatus{none}
  \notclaimed{This chapter does not claim open-cover compactness, selected subcovers, quotient real equality, countable choice, ambient RealUp completeness, Lean verification, or bridge checking.}
  \upgradepath{Obligation closure requires checked carrier admission, compact interval source exposure, finite-net handoff, tolerance routing, readback, RealUp sealing, transport, replay, provenance, and local naming for BEDC.Derived.CompactIntervalModulusUp.}
\end{closurestatus}
